% Auto-generated by scripts/make_tables.py -- do not hand-edit.
% Requires: \usepackage{booktabs} \usepackage{graphicx}
% The \resizebox around the tabular fits it to the text width;
% use it inside a full-width float (table*) in a two-column paper.
\begin{table}[!htbp]
\centering
\caption{Protocol span versus training gain, per task.}
\label{tab:span-vs-gain}
\footnotesize
\resizebox{\linewidth}{!}{%
\begin{tabular}{p{0.34\linewidth}ccc}
\toprule
 & SR $\times$4 & SR $\times$2 & Low-light \\
\midrule
Protocol span, P2-P8 (dB) & 3.889 & 3.143 & 2.250 \\
Protocol span, incl. best-of-60 P1 (dB) & 19.794 & 15.648 & 6.067 \\
Training gain under P3 (dB) & 0.782 & 0.863 & 10.736 \\
Span / training gain & 4.98 & 3.64 & 0.21 \\
Span incl. P1 / training gain & 25.33 & 18.13 & 0.57 \\
Seed sigma, full-image protocols (dB) & 0.0190-0.0361 & 0.0268-0.0938 & 0.0346-0.0806 \\
Seed sigma, P1 best-of-60 (dB) & 0.2364 & 0.5584 & 0.4148 \\
Seed sigma, P2 centre crop (dB) & 0.0363 & 0.0685 & 0.5303 \\
\bottomrule
\end{tabular}}
\vspace{2pt}
\begin{minipage}{0.98\linewidth}\footnotesize\raggedright
\textbf{(a)} Every quantity is same-task on both sides of every ratio. No cross-task pooled interval is given: the low-light task is a counterexample (ratio 0.21), so a pooled '5-25x' would conceal it.\\
\textbf{(b)} The low-light training gain (10.74 dB) is 4.8x its protocol span, so on that task the protocol does not dominate the conclusion. The claim is conditional: when headroom and protocol effect are of the same order, the protocol decides the conclusion.\\
\textbf{(c)} Seed sigma covers only the full-image protocols. P1 and P2 sample, so their sigma is a different quantity and is listed in its own two rows: P1 draws 60 crops per image and carries the largest spread, while P2 is a single fixed crop and behaves like the full-image protocols on the SR tasks.\\
\end{minipage}
\end{table}
